% 本文件由 paperswarm.tex_gate 从台账自动生成，请勿手工编辑。
% 每个数值均由证据台账注入：claim_id -> (artifact SHA-256, JSON Pointer)。
\subsection{Discussion}

The evidence presented in this section highlights the effectiveness of ridge regression over the minimum-norm ordinary least squares (OLS) solution in reducing mean test mean squared error (MSE). Specifically, the relative reduction in mean test MSE achieved by ridge regression over OLS is 50.94\%, indicating a significant improvement in model generalization. The mean test MSE of the ridge regression solution is 0.3041, with a standard deviation of 0.0623, suggesting that the model's performance is more consistent across different random initializations.

In contrast, the OLS solution exhibits higher variability, as evidenced by its mean test MSE of 0.6198 and a standard deviation of 0.1329. This variability is a critical factor in the overall performance of the model, as it can lead to unreliable predictions. The selected ridge penalty (0.1) was chosen to minimize the mean test MSE, demonstrating the importance of tuning the regularization parameter to achieve optimal performance.

The irreducible noise floor (0.1225) serves as a reference point for the inherent variability in the data, which is essential for understanding the true performance of the models. The number of random seeds (8) over which the experiments were averaged ensures robustness and reliability of the results. With 120 features and 90 training samples per split, the experiments were conducted on a balanced dataset, allowing for fair comparison between OLS and ridge regression.

These findings support the use of ridge regression in scenarios where model generalization is critical, and the risk of overfitting must be mitigated. The high variance observed in the OLS solution (0.1329) underscores the necessity of employing regularization techniques to improve model stability and performance. Future work could explore the impact of different regularization methods and hyperparameter settings on model performance, as well as the application of these techniques to more complex and high-dimensional datasets.
